\documentclass[11pt,a4paper]{article}
\usepackage[margin=27mm]{geometry}
\usepackage{amsmath,amssymb,amsthm,booktabs,hyperref}
\hypersetup{hidelinks}
\newtheorem{lemma}{Lemma}
\newtheorem{theorem}[lemma]{Theorem}
\newcommand{\im}{\operatorname{im}}
\newcommand{\Soc}{\operatorname{Soc}}
\newcommand{\Aut}{\operatorname{Aut}}
\title{Two ternary codes from flags and Fano subplanes}
\author{Private derivation and finite certificate: review pending}
\date{}
\begin{document}\maketitle

This note addresses the $120$-coordinate modules $N$ and $K$ in
L050, Theorem~1(b) and Proposition~3. The source is Leemans--Rodrigues,
\emph{Linear codes with complementary duals from some strongly regular
subgraphs of the McLaughlin graph}, Math. Commun.\ 21 (2016), 239--249,
particularly printed pages 240 and 247. Its construction of the residual
Steiner graph is retained. We supply an incidence operator producing the
$15$-dimensional code, followed by complete finite rank, module, weight and
permutation-group certificates. The $21$-dimensional graph code is a different
subspace; its rank or LCD property does not establish the assertions here.
The original uses constructive MeatAxe arguments and retained representation
theory. Our discovery also used MeatAxe, but the final module certificate
checks an explicit filtration and all linear intertwiner spaces instead.
No historical novelty or whole-paper recovery is asserted.

\section{The incidence operator}
Use the fixed $S(4,7,23)$ already constructed and checked by
\texttt{steiner\_construction.py}. Delete two named points. The $120$ heptads
avoiding them form a set $\mathcal H$. The $21$ heptads containing both
points, with those two removed, form the lines of a projective plane on
$21$ points. Its lines have size five; every two remaining points lie on
one line by the Steiner property. Two such lines meet in one point, since
the original heptads meet in one or three points and already share the
deleted pair.

For $H\in\mathcal H$, every two points of $H$ lie on a plane line meeting
$H$ in three points. Indeed the heptad defining that line shares at least
two points with $H$ and hence exactly three. These triples form the Fano
subplane on the seven points of $H$. Let $\mathcal F$ be the $105$ incident
point--line flags of the plane, and define a $105\times120$ matrix $R$ by
\[
 R_{(p,\ell),H}=1\quad\Longleftrightarrow\quad
 p\in H\ \hbox{ and }\ |\ell\cap H|=3.
\]
Thus each column records the $21$ flags of the corresponding Fano subplane.
All following code spaces and matrices are over $\mathbb F_3$.

\begin{lemma}\label{lem:square}
If $R$ is any matrix over a field with $RR^{\mathsf T}=0$, then
$Q=R^{\mathsf T}R$ is symmetric and $Q^2=0$. Its image is self-orthogonal.
If also $Rj=0$ and $j^{\mathsf T}j=0$, then $\im Q+Fj$ is self-orthogonal.
\end{lemma}
\begin{proof}
The matrix identities give symmetry and
$Q^2=R^{\mathsf T}(RR^{\mathsf T})R=0$. Hence
$\langle Qx,Qy\rangle=x^{\mathsf T}Q^2y=0$.
Moreover $Qj=R^{\mathsf T}Rj=0$, so $j$ is orthogonal to $\im Q$;
its own norm is zero by hypothesis.
\end{proof}

For the fixed family, the exact integer scalar products of incidence rows
are as follows. This table is checked for every pair of the $105$ flags by
the retained complete construction, not inferred from a sampled orbit.
\[
\begin{array}{l|r}
\text{relation between the two flags}&(RR^{\mathsf T})_{f,g}\\ \hline
f=g&24\\
\text{same point or same line, but distinct}&12\\
\text{distinct points and lines, with one cross incidence}&6\\
\text{distinct points and lines, with neither cross incidence}&3.
\end{array}
\]
Two distinct points and distinct lines cannot have both cross incidences.
Thus the table exhausts all cases and proves $RR^{\mathsf T}=0$ over
$\mathbb F_3$. Each row has sum $24$, so $Rj=0$, and the constant vector on
$120$ coordinates has norm zero. Set
\[
 K=\im(R^{\mathsf T}R),\qquad N=K+\mathbb F_3j.
\]
The exact row reductions give $\dim K=15$ and $\dim N=16$.
Consequently the sum defining $N$ is direct. Here the direct sum is an
internal decomposition of modules on the same $120$ coordinates, not a
concatenation of codes on disjoint coordinate sets. Lemma~\ref{lem:square} proves
self-orthogonality of both codes. These ranks and the full table, including
its underlying block lists, are obligations of the complete finite certificate.

\section{Recognizing the socle}
The literal permutations in \texttt{module\_data.json} preserve the residual Steiner graph and generate a primitive
permutation group $G\cong L_3(4)$ on $\mathcal H$, with point stabilizer
$L_3(2)$. The certificate checks primitivity, the isomorphisms and stabilizer, so this is
the specified degree-$120$ coset action in the source. It also checks
invariance of $K$ and $N$. The following elementary criterion makes the
remaining module calculation transparent.

\begin{lemma}\label{lem:socle}
Let $V$ be a finite-dimensional $FG$-module with an invariant filtration
$0=V_0<V_1<\cdots<V_s=V$, whose factors $D_i=V_i/V_{i-1}$ are simple.
Let $N\subseteq V$ be a sum of simple submodules. If every $G$-homomorphism
$D_i\to V$, for every $i$, has image in $N$, then $\Soc(V)=N$.
A matrix representation on a $d$-dimensional space is absolutely irreducible
if the algebra generated by its matrices is all $M_d(F)$.
\end{lemma}
\begin{proof}
If $S\subseteq V$ is simple, choose the least $i$ for which $S\subseteq V_i$.
The intersections $S\cap V_j$ are submodules of $S$. Minimality implies
$S\cap V_{i-1}=0$, so projection embeds $S$ into the simple factor $D_i$.
It is therefore an isomorphism onto that factor. Its inverse, followed by
$S\hookrightarrow V$, is a homomorphism $D_i\to V$; the hypothesis gives
$S\subseteq N$. This holds for every simple submodule, proving
$\Soc(V)\subseteq N$. The reverse inclusion follows from the hypothesis on $N$.
For the last assertion, any invariant subspace is invariant under the whole
matrix algebra. A nonzero vector can be sent to any chosen vector by an
endomorphism, so that subspace is the whole space. The same argument holds
after scalar extension because a basis for $M_d(F)$ remains a basis after
extension.
\end{proof}

Here $V=\mathbb F_3^{120}$. The supplied basis certifies invariant prefixes
of dimensions
\[
0,15,34,49,64,65,84,99,100,119,120.
\]
The ten factor dimensions and homomorphism-space dimensions are
\[
\begin{array}{c|rrrrrrrrrr}
i&1&2&3&4&5&6&7&8&9&10\\\hline
\dim D_i&15&19&15&15&1&19&15&1&19&1\\
\dim\operatorname{Hom}_G(D_i,V)&1&0&0&0&1&0&1&1&0&1.
\end{array}
\]
Triangularity of the generator matrices verifies invariance of all prefixes.
For each factor the matrix algebra has dimension $(\dim D_i)^2$, proving
simplicity by Lemma~\ref{lem:socle}. The first prefix is exactly $K$.
For each factor, the verifier forms \emph{all} equations
$XP_g=D_gX$ for the given generators and computes their full solution space.
Every row of every basis solution belongs to $N$. Linearity then establishes
the hypothesis for every homomorphism, not only a selected map.
The constant line is simple and $K$ is the first simple factor, so
Lemma~\ref{lem:socle} yields
\[
 \Soc(\mathbb F_3^{120})=N=K\oplus\mathbb F_3j.
\]
This calculation does not classify every irreducible module of $L_3(4)$.
The source's classification assertion concerning three nonisomorphic faithful
irreducible modules of dimension $15$ is retained from its reference [13], Jansen--Lux--Parker--Wilson, \emph{An Atlas of Brauer Characters} (1995).
Our certificate identifies a particular such simple module; its faithfulness
also follows because $G$ is simple and the representation is nontrivial.

\section{Weights and full permutation automorphisms}
Complete enumeration in the certified bases gives
\[
\begin{array}{c|r|r|r|r|r}
\text{code}&\text{number of words}&\text{dimension}&d&\text{words of weight }d&d^\perp\\\hline
K&14348907&15&48&1260&4\\
N&43046721&16&48&1260&4.
\end{array}
\]
The enumerator visits each of the $3^k$ coefficient vectors exactly once;
independence of the basis makes these distinct words. It records the full
weight distribution. The dual distributions are obtained from the classical
MacWilliams identity
\[
 A_j^\perp=3^{-k}\sum_{i=0}^{120}A_i
 \sum_{s=0}^j(-1)^s2^{j-s}\binom{i}{s}\binom{120-i}{j-s},
\]
with binomial coefficients outside their ranges taken as zero. In both
cases the dual coefficients at weights $1,2,3$ vanish. At weight $4$ they are
$2100$ for $K^\perp$ and $1260$ for $N^\perp$. Thus the dual dimensions are $105$
and $104$, respectively, and the stated dual minimum distances follow.
These are complete finite computations, not distance bounds inferred from
self-orthogonality.

The orbit under $G$ of the supplied weight $4$8 word contains $1260$ distinct
words and $630$ distinct supports. Since $K\subset N$ and both codes have
exactly $1260$ minimum words, this orbit is the complete minimum-word set of
both codes. Every coordinate-permutation automorphism of either code
therefore acts on its support-incidence structure, with $120$ coordinate
vertices and $630$ distinct support vertices.

Here is a short upper-bound certificate for the \emph{full} automorphism
group of that structure. Give the coordinate and support vertices two
different initial colours. Refine colours by old colour and the multiset
of neighbour colours. Successively individualize the coordinate vertices
$1,6,48,10,2,3$ in the literal certificate ordering. Their successive
refined colour classes have sizes
\[
120,21,2,4,2,2.
\]
After these choices all coordinate classes are singletons. Since the support
neighbourhoods are distinct, an automorphism fixing the coordinates also
fixes the supports. An automorphism fixing the previous choices must
preserve every refined colour class. The successive orbit--stabilizer
bounds therefore give order at most
$120\cdot21\cdot2\cdot4\cdot2\cdot2=80640$.

The explicit permutations supplied in the certificate generate a group
$A$ of order $80640$ and preserve each of $K,N$. Thus both code automorphism
groups attain the upper bound and equal $A$. The certificate checks
$A'=G$ and supplies two commuting involutions generating a complement of
order $4$, disjoint from $G$. Consequently
\[
 \Aut(K)=\Aut(N)=A\cong L_3(4):2^2.
\]
No graph-isomorphism program is needed for this upper bound. Such software
was useful in discovering the permutations, but the retained proof checks
literal permutations, ordinary group orders and the full refinement chain.

\paragraph{Limits.}
This note recovers only the stated socle and associated-code assertions of
L050, Theorem 1(b)/Proposition 3. It does not establish the other source codes,
minimum-dimension or global module-classification claims. It contains no
new Lean formalization. The complete finite construction, module data,
linear-algebra verifier and weight enumerator are required mathematical
inputs and must be retained and charged in full when this proof is used.
No numerical compression factor or historical priority claim is made.
\end{document}
